\setcounter{section}{6}








Es seien
\mathkor {} {X} {und} {Y} {}
\definitionsverweis {topologische Räume}{}{.}
Eine
\definitionsverweis {stetige Abbildung}{}{}
\maabbdisp {p} { Y } { X
} {}
heißt
\definitionswort {Überlagerung}{,}
wenn es eine
\definitionsverweis {offene Überdeckung}{}{}

\mavergleichskette
{\vergleichskette
{X
}
{ = }{\bigcup_{i \in I} U_i
}
{  }{
}
{    }{
}
{   }{
}
}
{}{}{}
und eine Familie
\definitionsverweis {diskreter}{}{}
topologischer Räume

\mathbed {F_i} {}
 {i \in I} {}
 {} {} {} {,}
derart gibt, dass
\mathl{p^{-1}(U_i)}{}
\definitionsverweis {homöomorph}{}{}
zu
\mathl{U_i \times F_i}{}
\zusatzklammer {versehen mit der
\definitionsverweis {Produkttopologie}{}{}} {} {}
ist, wobei die Homöomorphien mit den Abbildungen nach $U_i$ verträglich sind.







\inputaufgabe
{}
{





Zeige, dass die Abbildung
\maabbeledisp {} {\R } { S^1
} { t } { \left(  \cos  t , \sin  t  \right)
} {,}
eine
\definitionsverweis {Überlagerung}{}{}
ist.





}
{} {}

\inputaufgabe
{}
{





Zeige, dass die Abbildung
\maabbeledisp {} { {\mathbb C} } { {\mathbb C} ^{\times} = {\mathbb C} \setminus \{0\}
} { z } { \exp  z
} {,}
eine
\definitionsverweis {Überlagerung}{}{}
ist.





}
{} {}

\inputaufgabe
{}
{





Zeige, dass es bei einer
\definitionsverweis {Überlagerung}{}{}
\maabb {p} { Y } { X
} {}
zu jedem Punkt

\mavergleichskette
{\vergleichskette
{ x
}
{ \in }{ X
}
{  }{
}
{    }{
}
{   }{
}
}
{}{}{}
eine
\definitionsverweis {offene Umgebung}{}{}

\mavergleichskette
{\vergleichskette
{ x
}
{ \in }{ U
}
{ \subseteq }{ X
}
{    }{
}
{   }{
}
}
{}{}{}
und einen
\definitionsverweis {stetigen Schnitt}{}{}
\maabb {s} { U } { p^{-1}(U)
} {}
zu $p$ gibt.





}
{} {}

\inputaufgabe
{}
{





Es seien
\mathkor {} {F} {und} {H} {}
kommutative
\definitionsverweis {topologische Gruppen}{}{}
und es sei

\mavergleichskette
{\vergleichskette
{ G
}
{ = }{ F \times H
}
{  }{
}
{    }{
}
{   }{
}
}
{}{}{}
ihre
\definitionsverweis {Produktgruppe}{}{}
\zusatzklammer {versehen mit der
\definitionsverweis {Produkttopologie}{}{}} {} {}
und es sei

\mathdisp {0 \, \stackrel{  }{\longrightarrow} \,  F   \, \stackrel{  }{ \longrightarrow}   \,  G   \, \stackrel{  }{\longrightarrow} \,  H \,\stackrel{  } {\longrightarrow}  \, 0} {  }

die zugehörige kurze exakte Sequenz. Zeige, dass zu jedem
\definitionsverweis {topologischen Raum}{}{}
$X$ eine
\definitionsverweis {kurze exakte Garbensequenz}{}{}

\mathdisp {0 \, \stackrel{  }{\longrightarrow} \, C^0(-,F)  \, \stackrel{  }{ \longrightarrow}   \, C^0(-,G)  \, \stackrel{  }{\longrightarrow} \, C^0(-,H) \,\stackrel{  } {\longrightarrow}  \, 0} {  }

vorliegt, für die die globale Auswertung hinten stets surjektiv ist.





}
{} {}

\inputaufgabe
{}
{





Es sei
\maabb {\varphi} { X } { Y
} {}
eine
\definitionsverweis {stetige Abbildung}{}{,}

\mathl{Q \in Y}{} ein Punkt und ${ \mathcal  F   }$ eine
\definitionsverweis {Prägarbe}{}{}
auf $X$. Zeige, dass der
\definitionsverweis {Halm}{}{}
der
\definitionsverweis {vorgeschobenen Prägarbe}{}{}

\mathl{\varphi_* { \mathcal  F   }}{} im Punkt $Q$ gleich

\mavergleichskettedisphandlinks
{\vergleichskettedisphandlinks
{ \operatorname{colim}\,_{Q \in V } { \mathcal  F   } \left(  \varphi^{-1}(V)  \right)
}
{ =} { \operatorname{colim}\,_{ \left\{  U \subseteq X  \mid  \text{ es gibt } Q \in V \text{ mit } \varphi^{-1}(V) \subseteq U         \right\}  } { \mathcal  F   } \left(  U   \right)
}
{ } {
}
{ } {
}
{ } {
}
}
{}{}{}
ist.





}
{} {}

\inputaufgabe
{}
{





Es sei
\maabb {\varphi} { X } { Y
} {}
einer
\definitionsverweis {stetige Abbildung}{}{}
und ${ \mathcal  G   }$ eine
\definitionsverweis {Garbe}{}{}
auf $Y$. Zeige, dass der
\definitionsverweis {Halm}{}{}
der
\definitionsverweis {zurückgezogenen Garbe}{}{}
in einem Punkt

\mavergleichskette
{\vergleichskette
{P
}
{ \in }{ X
}
{  }{
}
{    }{
}
{   }{
}
}
{}{}{}
gleich dem Halm von ${ \mathcal  G   }$ in
\mathl{\varphi(P)}{} ist.





}
{} {}

\inputaufgabe
{}
{





Es sei $X$ eine Menge mit zwei Topologien
\mathkor {} {\tau_1} {und} {\tau_2} {}
derart, dass die Identität
\maabbdisp {\varphi} {X_1 = (X, \tau_1)} {X_2 =  (X, \tau_2)
} {}
\definitionsverweis {stetig}{}{}
ist, die erste Topologie ist also eine Verfeinerung der zweiten Topologie. Es sei
\mathl{{ \mathcal  F   }_1}{} eine
\definitionsverweis {Garbe}{}{}
auf $X_1$ und
\mathl{{ \mathcal  F   }_2}{} eine Garbe auf $X_2$. Bestimme $\varphi_* { \mathcal  F   }_1$ und $\varphi^{-1} { \mathcal  F   }_2$. Wie sieht es aus, wenn $\tau_1$ die
\definitionsverweis {diskrete Topologie}{}{}
und $\tau_2$ die
\definitionsverweis {Klumpentopologie}{}{}
ist.





}
{} {}

\inputaufgabe
{}
{





Es sei $X$ ein
\definitionsverweis {topologischer Raum}{}{}
und
\maabb {\varphi} { X } { \{ P \}
} {}
die konstante Abbildung. Es sei ${ \mathcal  F   }$ eine
\definitionsverweis {Garbe}{}{}
auf $X$. Bestimme $\varphi_* { \mathcal  F   }$.





}
{} {}

\inputaufgabe
{}
{





Es sei $X$ ein
\definitionsverweis {topologischer Raum}{}{}
und

\mavergleichskette
{\vergleichskette
{P
}
{ \in }{ X
}
{  }{
}
{    }{
}
{   }{
}
}
{}{}{}
ein Punkt mit der zugehörigen Abbildung
\maabb {i} {\{P\}} {X
} {.}
Es sei ${ \mathcal  F   }$ eine
\definitionsverweis {Garbe von kommutativen Gruppen}{}{}
auf $\{P\}$. Beschreibe die Garbe
\mathl{i_* { \mathcal  F   }}{} auf den offenen Mengen von $X$. Wie sehen die Halme von
\mathl{i_* { \mathcal  F   }}{} aus, wenn $P$ ein abgeschlossener Punkt ist?





}
{} {Vergleiche auch

Aufgabe 4.9
}

\inputaufgabe
{}
{





Es sei $X$ ein
\definitionsverweis {topologischer Raum}{}{}
und
\maabb {\varphi} { X } { \{ P \}
} {}
die konstante Abbildung. Es sei ${ \mathcal  G   }$ eine
\definitionsverweis {Garbe}{}{}
auf $\{P \}$. Bestimme $\varphi^{-1} { \mathcal  G   }$.





}
{} {}

\inputaufgabe
{}
{





Es sei
\maabb {\varphi} { X } { Y
} {}
eine
\definitionsverweis {stetige Abbildung}{}{}
zwischen den
\definitionsverweis {topologischen Räumen}{}{}
\mathkor {} {X} {und} {Y} {}
und es sei ${ \mathcal  F   }$ eine
\definitionsverweis {Garbe}{}{}
auf $X$. Zeige, dass es einen natürlichen
\definitionsverweis {Garbenmorphismen}{}{}
\maabbdisp {} { \varphi^{-1}(  \varphi_* { \mathcal  F   }) } { { \mathcal  F   }
} {}
auf $X$ gibt.





}
{} {}

\inputaufgabe
{}
{





Es sei
\maabb {\varphi} { X } { Y
} {}
eine
\definitionsverweis {stetige Abbildung}{}{}
zwischen den
\definitionsverweis {topologischen Räumen}{}{}
\mathkor {} {X} {und} {Y} {}
und es sei ${ \mathcal  G   }$ eine
\definitionsverweis {Garbe}{}{}
auf $Y$. Zeige, dass es einen natürlichen
\definitionsverweis {Garbenmorphismen}{}{}
\maabbdisp {} { { \mathcal  G   } } { \varphi_*(  \varphi^{-1} { \mathcal  G   })
} {}
auf $Y$ gibt.





}
{} {}

\inputaufgabe
{}
{





Es sei
\maabb {\varphi} { X } { Y
} {}
eine
\definitionsverweis {stetige Abbildung}{}{}
zwischen den
\definitionsverweis {topologischen Räumen}{}{}
\mathkor {} {X} {und} {Y} {.}
Es sei ${ \mathcal  F   }$ eine
\definitionsverweis {Garbe}{}{}
auf $X$ und ${ \mathcal  G   }$ eine Garbe auf $Y$. Zeige, dass es eine natürliche Bijektion zwischen den
\definitionsverweis {Garbenmorphismen}{}{}

\maabbdisp {\psi} { \varphi^{-1} { \mathcal  G   } } { { \mathcal  F   }
} {}
auf $X$ und den Garbenmorphismen
\maabbdisp {\theta} { { \mathcal  G   } } { \varphi_* { \mathcal  F   }
} {}
auf $Y$ gibt.





}
{} {}

\inputaufgabe
{}
{





Es seien
\mathl{L_1,L_2,M}{} Mengen und
\maabbdisp {p_1} {L_1 } { M
} {}
und
\maabbdisp {p_2} {L_2 } { M
} {}
Abbildungen. Wir definieren

\mavergleichskettedisp
{\vergleichskette
{L_1 \times_M L_2
}
{ \defeq} { \left\{ (x_1,x_2)  \mid   p_1(x_1) = p_2(x_2)         \right\}
}
{ \subseteq} { L_1 \times L_2
}
{ } {
}
{ } {
}
}
{}{}{.}
\aufzaehlungzwei {Zeige, dass es ein kommutatives Diagramm



\mathdisp {\begin{matrix} L_1 \times_M L_2  & \stackrel{  }{\longrightarrow} &  L_1  &  \\ \downarrow & & \downarrow  & \\  L_2  & \stackrel{  }{\longrightarrow} &  M & \! \! \! \!\!\!\!\!   \\ \end{matrix}} {  }




gibt.
} {Es sei $T$ eine weitere Menge und
\maabb {\psi_1} { T } { L_1
} {}
und
\maabb {\psi_2} { T } { L_2
} {}
Abbildungen mit

\mavergleichskettedisp
{\vergleichskette
{ p_1 \circ \psi_1
}
{ =} { p_2 \circ \psi_2
}
{ } {
}
{ } {
}
{ } {}
}
{}{}{.}
Zeige, dass es eine eindeutig bestimmte Abbildung
\maabbdisp {\psi} { T } { L_1 \times_M L_2
} {}
derart gibt, dass die Projektionen auf
\mathkor {} {L_1} {bzw.} {L_2} {}
mit $\psi_1,\psi_2$ übereinstimmen.
}





}
{} {}

\inputaufgabe
{}
{





Es seien
\mathl{L_1,L_2,M}{}
\definitionsverweis {topologische Räume}{}{}
und
\maabbdisp {p_1} {L_1 } { M
} {}
und
\maabbdisp {p_2} {L_2 } { M
} {}
\definitionsverweis {stetige Abbildungen}{}{.}
Wir definieren

\mavergleichskettedisp
{\vergleichskette
{L_1 \times_M L_2
}
{ \defeq} { \left\{ (x_1,x_2)  \mid   p_1(x_1) = p_2(x_2)         \right\}
}
{ \subseteq} { L_1 \times L_2
}
{ } {
}
{ } {
}
}
{}{}{}
mit der
\definitionsverweis {induzierten Topologie}{}{.}
\aufzaehlungzwei {Zeige, dass es ein kommutatives Diagramm



\mathdisp {\begin{matrix} L_1 \times_M L_2  & \stackrel{  }{\longrightarrow} &  L_1  &  \\ \downarrow & & \downarrow  & \\  L_2  & \stackrel{  }{\longrightarrow} &  M & \! \! \! \!\!\!\!\!   \\ \end{matrix}} {  }




mit stetigen Abbildungen gibt.
} {Es sei $T$ ein weiterer topologischer Raum und
\maabb {\psi_1} { T } { L_1
} {}
und
\maabb {\psi_2} { T } { L_2
} {}
stetige Abbildungen mit

\mavergleichskettedisp
{\vergleichskette
{ p_1 \circ \psi_1
}
{ =} { p_2 \circ \psi_2
}
{ } {
}
{ } {
}
{ } {}
}
{}{}{.}
Zeige, dass es eine eindeutig bestimmte stetige Abbildung
\maabbdisp {\psi} { T } { L_1 \times_M L_2
} {}
derart gibt, dass die Projektionen auf
\mathkor {} {L_1} {bzw.} {L_2} {}
mit $\psi_1,\psi_2$ übereinstimmen.
}





}
{} {}

\inputaufgabe
{}
{





Es seien
\mathkor {} {X} {und} {Y} {}
\definitionsverweis {topologische Räume}{}{}
und
\maabb {\varphi} { Y } { X
} {}
eine
\definitionsverweis {stetige Abbildung}{}{.}
Es sei
\maabb {p} { V } { X
} {}
ein
\definitionsverweis {Vektorbündel}{}{}
über $X$. Zeige, dass
\mathl{Y \times_X V}{}
\zusatzklammer {siehe

Aufgabe 6.15
} {} {}
ein Vektorbündel über $Y$ ist.





}
{} {}

\inputaufgabe
{}
{





Es sei $X$ ein
\definitionsverweis {topologischer Raum}{}{}
und seien
\maabb {p} { V } { X
} {}
und
\maabb {q} { W } { X
} {}
\definitionsverweis {Vektorbündel}{}{}
über $X$. Zeige, dass
\mathl{V \times_X W}{}
\zusatzklammer {siehe

Aufgabe 6.15
} {} {}
ein Vektorbündel über $X$ ist, das mit der
\definitionsverweis {direkten Summe von Vektorbündeln}{}{}
über $X$ übereinstimmt.





}
{} {}

\inputaufgabe
{}
{





Es seien
\mathkor {} {X,Y} {und} {Z} {}
\definitionsverweis {topologische Räume}{}{}
und seien
\maabb {\varphi} { Y } { X
} {}
und
\maabb {p} { Z } { X
} {}
\definitionsverweis {stetige Abbildungen}{}{.}
Es sei
\maabbdisp {p_Y} {Y \times_XZ } { Y
} {.}
Zeige, dass ein
\definitionsverweis {stetiger Schnitt}{}{}
\maabb {s} { Y } { Y \times_XZ
} {}
das gleiche ist wie eine stetige Abbildung
\maabb {t} { Y } { Z
} {}
mit

\mavergleichskette
{\vergleichskette
{ p \circ t
}
{ = }{ \varphi
}
{  }{
}
{    }{
}
{   }{
}
}
{}{}{.}





}
{} {}

\inputaufgabe
{}
{





Es seien
\mathkor {} {X,Y} {und} {Z} {}
\definitionsverweis {topologische Räume}{}{}
und seien
\maabb {\varphi} { Y } { X
} {}
und
\maabb {p} { Z } { X
} {}
\definitionsverweis {stetige Abbildungen}{}{.}
Es sei
\maabb {p_Y} {Y \times_XZ } { Y
} {.}
Es sei ${ \mathcal  G   }$ die
\definitionsverweis {Garbe der stetigen Schnitte}{}{}
zu $p$ auf $X$. Zeige, dass der
\definitionsverweis {Rückzug}{}{}
$\varphi^*{ \mathcal  G   }$ mit der Garbe der Schnitte zu $p_Y$ übereinstimmt.





}
{} {}
